\section{科研基础与研究经历}

我目前为浙江理工大学物理学专业硕士研究生，预计于2027年6月毕业。硕士阶段主要围绕低维功能材料、多铁耦合及相关量子物性开展第一性原理与多尺度模拟研究。

在科研训练过程中，我逐渐形成了较为明确的研究方式：从具体物理问题出发选择方法，通过计算、模型和程序建立可验证的证据链，并进一步分析结构、电子态和微观相互作用之间的联系。相比单纯掌握某一种软件，我更重视问题发现、方法设计和结果验证能力。

目前，我较为熟悉VASP、Wannier90/WannierTools、TB2J和Spirit等研究工具，能够开展结构优化、电子结构、自旋轨道耦合、铁电极化与翻转路径、界面结构、磁交换作用、磁各向异性、DMI、Berry曲率及拓扑性质等计算，并根据研究问题建立有效模型和开展多尺度数值模拟。

在二维多铁材料TcIrGe$_2$Se$_6$的研究中，我从结构稳定性和铁电翻转出发，进一步分析磁基态、谷自由度、Berry曲率及微观交换机制，并利用Wannier紧束缚模型追踪Tc--Se--Ir--Se--Tc长程超超交换通道，将局域结构变化、轨道杂化与磁交换机制联系起来。该工作形成了“结构—极化—电子结构—微观相互作用—有效模型—功能响应”的完整研究链条，相关论文已被\textit{Physical Review B}接收。

此外，我还开展了成分调控和二维铁电/磁性异质结构研究，重点关注极化翻转、界面构型、层间耦合和应变对材料功能性质的影响。这些研究使我的兴趣逐渐从单一物性预测转向铁电相变、亚稳态、界面结构及真实材料可实现性等更具普适性的问题。

\section{创新能力与方法开发}

我认为自己的一个主要科研特点，是在现有分析方法不足以回答问题时，愿意进一步开发新的分析工具。

在研究TcIrGe$_2$Se$_6$长程交换机制时，仅依赖常规能带和投影态密度难以定量判断具体超超交换路径。为此，我基于Wannier模型自行编写Python程序，解析\texttt{wannier90\_hr.dat}、Wannier中心和晶体结构信息，对不同原子、轨道及晶格平移之间的跃迁矩阵元进行筛选和统计，从而定量提取特定轨道路径的杂化强度。这一过程使原本依赖定性判断的交换机制分析转化为可计算、可比较、可复核的定量证据。对我而言，这类从科学问题出发主动建立分析方法的过程，是科研创新能力的重要体现。

我也在尝试将机器学习与物理模型结合。在交错磁拓扑研究中，我将Sobol参数采样、Chern数计算、分层机器学习和局部相界搜索结合，用代理模型缩小高维参数空间中的候选区域，再返回严格的物理计算进行验证。整个项目最终按照不同物理模块整理为5个代码单元、34个Python源文件，并区分训练样本、外部验证样本、代表参数和仍需进一步验证的问题。

这项经历使我认识到，机器学习在材料研究中的价值并不是取代物理模型，而是帮助研究者更有效地搜索参数空间、识别规律，并将计算资源集中到真正值得验证的问题上。

\section{编程能力与科研工作能力}

编程是我开展科研的重要工具。我能够使用Python和Shell/Linux进行批量计算、任务管理、参数扫描、数据提取、统计分析及结果可视化，并能够根据课题需要自行搭建自动化流程。

长期计算项目使我逐渐形成了程序化、模块化和可追溯的工作习惯。对于复杂课题，我通常会将任务拆分为结构、电子、极化、模型、后处理和论文写作等模块，并记录计算参数、图件来源、异常结果及后续验证任务。这种方式使多个长期项目能够并行推进，也能够在经历多轮修改后快速恢复研究上下文。

我同时重视计算结果的可靠性。针对U值、SOC设置、$k$点密度、Wannier拟合窗口、有限尺寸模型和边界条件等可能影响结论的因素，我会主动设计收敛测试和交叉验证。当不同方法得到不一致结果、模型出现多个亚稳态或计算趋势偏离预期时，我通常会重新检查结构、参数和物理假设，并通过补充计算或修改程序定位问题。

我享受“发现异常—提出假设—设计验证—形成解释”的科研过程。相比得到一个符合预期的结果，我更关注该结果是否具有足够可靠的物理依据。目前，我已有1篇\textit{Physical Review B}论文接收、1篇相关论文在审，另有二维异质结构和机器学习方向工作正在推进。论文撰写和修改经历也使我更加重视计算数据、图表、方法和文字表述之间的一致性。

\section{博士阶段研究计划}

博士阶段，我希望将研究重点逐步转向铁电/反铁电材料、外延薄膜、异质界面及极化调控机制，尤其关注以下问题：
\begin{enumerate}[label=\arabic*.,leftmargin=2.2em,itemsep=0.05em,topsep=0.15em]
  \item 极化翻转过程中的原子尺度结构演化；
  \item 铁电相、反铁电相及亚稳态之间的竞争；
  \item 应变、界面、缺陷和晶向对极化稳定性及翻转路径的影响；
  \item 复杂极化畴和拓扑极化结构的形成及外场调控。
\end{enumerate}
这些问题与我目前掌握的第一性原理、NEB、界面建模、编程和机器学习方法具有较好的连续性。

与此同时，我希望在博士阶段系统补充实验能力，学习外延薄膜制备、铁电性能测试、PFM、高分辨STEM及同步辐射结构表征等实验技术。我的目标并不是简单地从理论计算“转向实验”，而是逐步建立理论与实验能够相互反馈的研究模式：

\begin{center}
\sffamily\small\bfseries\color{StatementTeal}
理论预测与机制分析 $\rightarrow$ 材料与结构设计 $\rightarrow$ 薄膜/异质结构制备\\[-0.1em]
$\rightarrow$ 电学及原子尺度表征 $\rightarrow$ 实验反馈修正模型 $\rightarrow$ 再设计与再验证
\end{center}

在这一过程中，我也希望发挥已有的编程和机器学习优势。例如，通过第一性原理数据构建机器学习原子势，研究传统DFT难以覆盖的大尺度极化翻转、畴壁运动及有限温度结构演化；或通过数据驱动方法分析成分、应变、取向及界面结构与铁电性能之间的关系。我的目标是使机器学习真正服务于材料物理问题，而不是成为脱离物理机制的独立工具。

\section{个人特点与发展目标}

硕士阶段的科研经历使我越来越确认，自己真正感兴趣的是发现问题、建立假设、设计验证并最终理解微观机制的过程。我具有较强的自主学习和方法迁移意愿：面对新的科学问题，我愿意学习新的理论、程序和实验方法；面对异常结果，我愿意继续追问并重新验证；面对长期复杂项目，我能够通过程序化、模块化和文档化的方式持续推进。

博士阶段，我希望进一步提升两个方面的能力：一是从目前能够独立完成计算课题，进一步成长为能够独立发现并凝练科学问题的研究者；二是突破目前以计算为主的研究边界，逐渐形成理论计算、数据方法和实验研究相结合的能力。

长期而言，我希望继续在高校、科研院所或高水平科研平台从事科研工作，并围绕铁电/多铁材料、低维功能材料、界面工程及极化调控形成具有连续性的研究方向。我希望未来形成的个人研究特色，不是局限于某一种软件或实验技术，而是能够针对具体科学问题选择合适的方法，并逐步贯通“材料设计—理论预测—程序与数据方法—材料制备—原子尺度表征—功能测试—物理机制分析”。

我希望通过博士阶段更加系统的科研训练，从目前以计算研究为主的硕士生，成长为具有独立创新能力、完整科研执行能力以及理论—实验交叉研究能力的科研工作者。